\chapter{基于规则的知识获取}
\label{chap:rule-based-knowledge-acquisition}

假设我们要从中文新闻中找出任职变动，并记录人员与组织之间的任职关系。一条规则只要看到“任命”就输出“任职变动”，便会同时命中“海岬科技任命林舟为首席架构师”“海岬科技未任命林舟”和“关于任命林舟的传言已被否认”。三句话都有这个词，只有第一句直接支持目标事件。程序找到词语，只完成了字面检查；还要判断规则是否应当触发，以及触发后的结论是否符合任务定义。

要让这类判断能够重复执行，我们可以把检查条件、查询已有知识的办法或答案必须满足的约束写成规则。\term{基于规则的知识获取}（rule-based knowledge acquisition）就是按这些明确依据，为样本产生任务标签、候选范围或关系信息。所得内容仍是候选知识，需要说明它凭什么产生。专家编写、从数据归纳和由模型生成回答的是\emph{规则怎样形成}；检查样本模式、关联已有知识和求解约束回答的是\emph{规则依据什么产生知识}。例如，模型可以提出一条查询知识库事实的规则，专家也可以写出只比较数值阈值的规则。提出规则的来源与规则使用的判断依据应分开记录。

在上述新闻任务中，没有出现任命、无法分析否定语境和输入不属于新闻，是三种不同情况。为了把它们写入第~\ref{chap:knowledge-acquisition-overview}章中含六个字段的知识记录，我们需要分别计算适用前提与触发条件，再据此确定有无输出。设样本为$x$，规则版本为$v$，适用前提为$A_v(x)$，触发条件为$C_v(x)$，输出构造为$O_v(x)$。本章用$\{\mathsf{T},\mathsf{F},\mathsf{U}\}$分别表示条件成立、不成立和无法确定。\term{强Kleene合取}（strong Kleene conjunction）规定：任一项为$\mathsf{F}$则合取为$\mathsf{F}$，全为$\mathsf{T}$才得$\mathsf{T}$，其余为$\mathsf{U}$。这使已确认不成立的条件可以决定不触发，而不会被另一个未知条件遮住。令$D_v(x)=A_v(x)\land C_v(x)$，一次成功完成条件检查后的结果分支为
\begin{equation}
  b_v(x)=
  \begin{cases}
    \text{范围外}, & A_v(x)=\mathsf{F},\\
    \text{触发}, & D_v(x)=\mathsf{T},\\
    \text{未触发}, & A_v(x)\ne\mathsf{F}\ \text{且}\ D_v(x)=\mathsf{F},\\
    \text{无法确定}, & D_v(x)=\mathsf{U}.
  \end{cases}
  \label{eq:rule-based-knowledge-acquisition-interface}
\end{equation}
式中的$b_v(x)$回答“这次检查走到哪个分支”，并不代替完整知识记录。适用前提$A_v(x)=\mathsf{F}$时，对象记为范围外，语义状态为“不适用”，来源动作为空。其余分支把“触发”“未触发”或“无法确定”保存在取得条件$s_{\mathrm{cond}}$的具名触发结果中。对于条件满足后即可构造任务内容的规则，触发时作答并给出$O_v(x)$，语义状态为已给出；未触发时没有来源动作，任务内容仍未知；合取无法确定时，以弃权动作保留未知内容。若输出还依赖后续求解，触发只表示进入求解；求解结果与最终来源动作仍需另记，不能把开始求解当作已经作答。

上述分支都以条件检查完成为前提，所以这一步的执行状态均为成功，答案却未必已经核验。后续求解是否完成，要保存其执行状态；程序中断应记执行失败，而不是主动弃权。范围外、未触发和弃权也都没有证明样本属于负类；只有规则明确检查了负类定义并输出负标签时，所得记录才能作为负例。

确定了何时作答，还要约定规则返回什么。输出$O_v(x)$可以是“任职变动”这一候选标签，可以是$\{\text{管理岗位},\text{技术岗位}\}$这样的候选范围，也可以是“两个实体不得取同一身份”这样的关系约束。复核这些内容所需的依据不同：匹配规则要留下命中的输入位置，查询规则要留下外部事实，推理规则要留下推理步骤；未能判断时则要保存原因。表~\ref{tab:rule-based-knowledge-acquisition-record}沿用任命例子，把这些记录要求放在一起。

\begin{table*}[htbp]
  \centering
  \begingroup
  \setlength{\tabcolsep}{4pt}
  \setlength{\parindent}{0pt}
  \setlength{\RaggedRightParindent}{0pt}
  \renewcommand{\arraystretch}{1.12}
  \begin{tabularx}{\linewidth}{>{\raggedright\arraybackslash}p{19mm}>{\raggedright\arraybackslash}p{32mm}>{\raggedright\arraybackslash}p{49mm}>{\raggedright\arraybackslash}X}
    \toprule
    字段 & 作用 & 教学取值 & 缺少时无法判断的问题 \\
    \midrule
    输入对象 & 确定规则检查什么 & 一句中文新闻及其中的实体提及 & 规则究竟作用于整篇文档、句子还是实体对 \\
    适用前提 & 限定允许使用规则的范围 & 中文新闻语体，事件时间符合任务口径 & 对象在范围外，还是缺少判断范围所需的证据 \\
    触发条件 & 说明何时产生输出 & “任命”连接组织与人员，且不在否定或被否认的引语中 & 命中由哪个条件造成，新增例外改变了什么 \\
    输出 & 说明产生何种候选知识 & 候选关系$(\text{林舟},\text{任职于},\text{海岬科技})$ & 返回的是类别、候选集合还是关系约束 \\
    无输出状态 & 区分范围、触发结果与动作 & 条件检查：范围外记不适用；未触发动作为\textemdash；合取无法确定时弃权 & 无任务答案是因为范围外、未触发，还是依据不足 \\
    命中证据 & 保留从输入到输出的路径 & 关键词跨度、两个实体标识及局部依存路径 & 程序成功是否等于样本语义支持结论 \\
    适用范围及版本 & 支持更新与追溯 & 新闻语体；规则\texttt{appointment-v1}；知识库2026-09版 & 词表、解析器或外部知识更新后哪些旧结果需要重算 \\
    \bottomrule
  \end{tabularx}
  \endgroup
  \caption{一条知识获取规则需要记录的内容。教学规则可以返回关系候选，也可以按任务改为类别或候选集合；字段的作用是把语义状态、来源动作、执行状态、知识依据和版本范围分开。}
  \label{tab:rule-based-knowledge-acquisition-record}
\end{table*}

沿着任命例子，我们可以分别检查句子中的词语与角色、查明“林舟”对应哪个实体，或用岗位关系限制候选答案。下面三节据此按知识产生的主要依据展开：模式匹配检查样本中的可观察条件，知识关联使用已有实体、概念或事实，约束推理限制候选答案之间的关系。一条规则可以同时使用这些办法，但仍要留下各步的依据，才能定位错误发生在匹配、关联还是求解环节。

% 标题、入口段与完整Snorkel边注共同检查剩余空间。
\Needspace{85mm}
\section{基于模式匹配的规则利用}
\label{sec:rule-based-knowledge-acquisition-pattern}

对一句待标注新闻，首先可用的材料就是句中出现的词语、位置和数值。我们可以检查这些材料是否满足预先写出的条件，条件成立时给出候选输出；这种直接检查样本的办法称为\term{模式匹配}（pattern matching）。Snorkel中的标注函数提供了一种可重复执行的实现：函数依据模式、启发式规则或外部资源为数据点给出标签，也可以在依据不足时弃权\scite{ratner2017}

计算模式之前，通常先把样本转换为分词结果、字段值或其他可检查的表示，再保存命中的跨度、数值和中间条件。这样就能区分三种障碍：字段缺失使条件无法确定，条件算出假值使规则未触发，条件虽成立却不支持目标语义使输出缺少依据。前两种影响触发分支，第三种需要检查已产生的候选；仅凭“没有可用答案”不能把它们归为执行失败。后面先看符号与数值两种基本条件，再看条件如何组合；它们都依赖输入格式和前处理，变化后需要重新检查。

\subsection{基于符号模式的规则利用}
\label{sec:rule-based-knowledge-acquisition-symbolic-pattern}

一句新闻中的“任命”是可观察的符号线索，但只看这个词还不能确定谁任命了谁。关键词检查词语是否出现；正则表达式可以进一步限制字符的位置与相邻形式；句式模板或局部结构模式则检查若干成分是否以规定关系出现。由文本扩展到日志和结构化记录时，检查对象可以换成字段或事件符号，判断依据仍是当前样本中的符号结构，而不是外部事实。

把任命规则的输入限定为已经切分的句子后，我们检查句内谓词及其参与者：任命者应是组织，被任命者应是人员。触发条件还要求谓词不受否定词支配，也不只出现在随后被否认的引语中。在声明范围内，全部条件成立才输出任职变动候选及关系三元组；存在假值就不触发，没有假值而仍有未知项才弃权。对“海岬科技任命林舟为首席架构师”，规则记录“任命”的字符跨度、海岬科技的组织跨度、林舟的人员跨度以及三者的局部结构。复核者由此可以回到输入，检查这些角色是否确实支持输出。

章首三条教学文本都能通过关键词检查，差别在任命是否被肯定地陈述。加入“任命谓词不受否定词支配”的条件，可以排除直接否定；被引用的任命传言随后遭到否认时，局部词语却仍可能呈肯定形式。规则因而还要检查引语边界和报道者是否否认了引语内容，无法可靠分析的条件记为未知。图~\ref{fig:rule-based-knowledge-acquisition-pattern-rule}把各项结果分别保存，再统一合取；我们应观察未知项与假值怎样共同决定分支。增加条件收窄了触发范围，但这些条件本身仍可能出错，不能据条件数量断言输出更准确。

\begin{figure*}[htbp]
  \centering
  \begin{tikzpicture}[
  x=1cm,
  y=1cm,
  font=\figurefont\footnotesize,
  panel/.style={draw=BookRule,line width=0.8pt,fill=white},
  text box/.style={draw=BookMuted,line width=0.8pt,fill=BookPaper,minimum height=10mm,text width=43mm,align=left,inner sep=2mm},
  condition/.style={draw=FigureModel,line width=1pt,fill=FigureModel!10,minimum height=10mm,text width=38mm,align=center,inner sep=1.5mm},
  output/.style={draw=FigureOutput,line width=1pt,fill=FigureOutput!7,minimum height=12mm,text width=39mm,align=center,inner sep=1.5mm},
  combine/.style={condition,draw=BookInk,fill=BookPaper},
  main/.style={knowledge arrow,draw=BookInk,line width=1.1pt,rounded corners=1mm}
]
  \fill[white] (0,0) rectangle (16.6,8);
  \draw[panel] (0,0) rectangle (5.15,8);
  \draw[panel] (5.4,0) rectangle (11.05,8);
  \draw[panel] (11.3,0) rectangle (16.6,8);

  \node[anchor=north west,font=\figurefont\bfseries] at (0.2,7.8) {(a) 相同字面命中};
  \node[text box,anchor=north west] at (0.25,7.1) {肯定：海岬科技\textbf{任命}林舟。\\\textcolor{BookTeal}{字面命中}};
  \node[text box,anchor=north west] at (0.25,4.95) {否定：海岬科技\textbf{未任命}林舟。\\\textcolor{BookTeal}{字面命中}};
  \node[text box,anchor=north west] at (0.25,2.8) {引用后否认：“\textbf{任命}林舟”的传言被否认。\\\textcolor{BookTeal}{字面命中}};
  \node[anchor=south,text=BookMuted,font=\figurefont\scriptsize,align=center] at (2.575,0.25) {匹配成功只说明找到了跨度};

  \node[anchor=north west,font=\figurefont\bfseries] at (5.6,7.8) {(b) 条件记录后统一合取};
  \node[condition] (scope) at (8.2,6.65) {$A_v$：适用前提\\记录$\mathsf{T}/\mathsf{F}/\mathsf{U}$};
  \node[condition] (keyword) at (8.2,5.35) {$K$：出现目标谓词\\记录$\mathsf{T}/\mathsf{F}/\mathsf{U}$};
  \node[condition] (roles) at (8.2,4.05) {$E$：组织与人员角色匹配\\记录$\mathsf{T}/\mathsf{F}/\mathsf{U}$};
  \node[condition] (context) at (8.2,2.75) {$N$：无否定或引语否认\\记录$\mathsf{T}/\mathsf{F}/\mathsf{U}$};
  \node[combine] (conjunction) at (8.2,1.0) {保留$A_v$并计算\\$D_v=A_v\land K\land E\land N$};
  \foreach \condition in {scope,keyword,roles,context} {
    \draw[main,-] (\condition.east) -- (10.7,0 |- \condition.east);
    \fill[BookInk] (10.7,0 |- \condition.east) circle (1pt);
  }
  \draw[main] (10.7,6.65) -- (10.7,1.95) -| (conjunction.north);

  \node[anchor=north west,font=\figurefont\bfseries] at (11.5,7.8) {(c) 分支与来源动作};
  \node[output] (triggered) at (14.1,6.65) {$D_v=\mathsf{T}$：触发、作答\\给出类别、范围或关系候选};
  \node[output,draw=BookMuted,fill=BookPaper] (untriggered) at (14.1,4.9) {$A_v\ne\mathsf{F},\ D_v=\mathsf{F}$\\未触发；动作：\textemdash\\任务内容未知};
  \node[output,draw=BookAmber,fill=BookAmber!8] (abstain) at (14.1,3.15) {$D_v=\mathsf{U}$：无法确定\\动作：弃权；任务内容未知};
  \node[output,draw=BookMuted,fill=white] (outside) at (14.1,1.4) {$A_v=\mathsf{F}$：范围外\\语义：不适用；动作：\textemdash};
  \coordinate (output-split) at (11.55,1.0);
  \draw[main,-] (conjunction.east) -- (output-split);
  \draw[main,-] (output-split) -- (11.55,6.65);
  \fill[BookInk] (output-split) circle (1pt);
  \foreach \result in {triggered,untriggered,abstain,outside} {
    \draw[main] (11.55,0 |- \result.west) -- (\result.west);
  }
\end{tikzpicture}

  \caption{从字面命中到带例外的组合规则。左面板中相同谓词出现在肯定、否定和被否认的引语里；中面板分别记录适用前提$A_v$、谓词条件$K$、角色条件$E$和非否定条件$N$，再计算$D_v=A_v\land K\land E\land N$，并列记录不要求并行执行。右面板按式~\eqref{eq:rule-based-knowledge-acquisition-interface}区分范围外、触发、未触发与无法确定；\textemdash{}表示来源动作为空。所有分支均以执行成功为前提。遇到未知仍须检查可独立求值的条件，例如$A_v=\mathsf{T}$且$(K,E,N)=(\mathsf{T},\mathsf{U},\mathsf{F})$时应未触发，而非弃权。}
  \label{fig:rule-based-knowledge-acquisition-pattern-rule}
\end{figure*}

上述角色与否定检查需要比关键词更多的前处理。字面关键词依赖规范化与切分，容易受多义词影响；带位置条件的正则表达式限制局部顺序，仍难处理跨句否定；句式模板检查实体角色和依存路径，又依赖实体识别与句法分析。规则使用的结构越多，需要核验的中间结果就越多。维护时不仅要修改规则文本，还要检查词形、拼写、分词方式、领域句式和解析器版本的变化。

怎样得到这些符号条件，是与执行相连的另一个问题。专家可以直接写出条件，我们也可以从语料中的重复结构归纳模板，或让模型提出候选。候选无论来自哪里，都要回到原始样本检查支持情形和反例，定位误触发原因，再调整条件或适用范围。触发位置与解析结果一并保留，才能判断是规则把正确结构解释错了，还是前处理没有得到正确结构。

即使符号检查执行正确，目标知识仍可能没有成立。若任务只标注已经生效的任职关系，准确匹配到尚未生效的任命公告也不能直接把它当作事实；反过来，解析依据不足使规则弃权时，样本实际仍可能属于目标类。前一种输出不符合任务定义，后一种未覆盖本来需要作答的样本。第~\ref{chap:knowledge-quality-assessment-correction}章将用统一的质量检查处理这些问题。

\subsection{基于数值条件的规则利用}
\label{sec:rule-based-knowledge-acquisition-numeric-pattern}

设备是否过热，要检查的是读数而不是新闻中的词语。对于测量值、计数、比例或距离，规则先取得可比较的数值，再判断它是否落在指定阈值或区间内；这种办法称为\term{数值条件规则}（numeric-condition rule）。比较计算没有算错，只能说明它忠实于输入与公式，不能说明阈值适用于当前设备或任务。

为看清边界，假设教学任务规定：设备读数不低于$80\,^\circ\mathrm{C}$时输出“高温候选”，在$[78,80)\,^\circ\mathrm{C}$内因测量误差影响而弃权，低于$78\,^\circ\mathrm{C}$时不触发这条规则。输入同时应有单位、采集时间和传感器状态。在其余前提满足时，$82\,^\circ\mathrm{C}$触发，$79\,^\circ\mathrm{C}$弃权；缺少单位而无法完成比较时也弃权。$176\,^\circ\mathrm{F}$换算为$80\,^\circ\mathrm{C}$，恰好达到包含端点的下界，所以触发。如果把数值$176$直接当作摄氏读数，比较仍能运行，得到的候选却依据了错误输入。

温度例子说明，数值规则不只是写一个比较符号。规则需要计算特征，统一单位或基准，执行严格或非严格比较，再把比较结果转换为触发、未触发或无法确定。使用多个区间时，要明确端点归属及是否重叠，并规定缺失、溢出和不可比较值怎样处理。不确定区间应来自测量误差、舍入范围或其他证据限制；误差来源未知时可以保守弃权，但不能把任意缓冲带称为经过验证的置信区间。

例中的阈值由教学任务定义给出，因此它规定的是标签口径；口径改变时，受影响的记录需要重新标注。若阈值来自领域标准，要保存适用环境和标准版本；若它由少量样本调出，则要保存开发数据及调整过程。经验阈值还需用未参与调参的数据独立检查，才能避免把开发样本上的吻合当作新样本上的依据。群体构成、采集设备或时间分布变化后，要重新判断原阈值是否适用于现在的对象。

阈值数值固定以后，取得候选仍需要计算特征、查询基准、换算单位并处理缺失值。如果不同地区、设备或人群使用不同基准，比较前还要选择与输入对应的基准。选错基准、算错特征和把范围外样本交给规则，会在不同环节产生错误；记录这些环节的输入与版本，维护者才能区分原因。监测输入变化也属于持续成本。

\subsection{基于组合条件的规则利用}
\label{sec:rule-based-knowledge-acquisition-composite-pattern}

任命规则只有同时确认谓词、实体角色和非否定语境，才有依据输出候选。我们把“出现任命谓词”“实体角色匹配”“没有否定或被否认引语”分别记为条件$K$、$E$和$N$，用$C_v(x)=K\land E\land N$组合它们。各项仍取章首规定的三值，并按强Kleene合取计算；实体识别器没有输出时，角色条件可能无法确定，不能直接把它记为实体类型不符合。

若同一任务有特殊情形和一般情形，我们还要规定先检查哪条规则、一个结果是否结束后续检查。\term{规则列表}（rule list）按优先级依次检查；本节采用保守策略，逐项条件检查确定未触发或范围外时才继续，触发时作答并停止，无法确定时弃权并停止。例如，特殊规则的$D_v(x)=\mathsf{U}$，即使一般规则能够触发，列表仍不采用一般答案，因为尚未排除特殊情形。默认分支只有在列表适用前提已确认、各项均已确定未触发或范围外，且任务明定这些结果足以判负时，才能输出负类；否则保留未触发且动作为空。\term{规则树}（rule tree）沿条件分支前进，本节同样规定未知时走专门的弃权出口，不把未知当作“否”。这两种未知处理是本节选定的控制策略，不是强Kleene合取本身的规定。互不设优先级的\term{规则集合}（rule set）则分别执行各条规则，保留全部原始输出，包括各条规则的未知结果。

逐项检查不一定要执行所有计算。只要某项假值已经确定合取为假，短路执行就可以跳过不影响结果的后续检查。例如，适用前提已确认、关键词为假时，无须继续运行昂贵的句法解析；关键词为真、角色未知时却还不能弃权，应继续检查能够独立求值的否定条件。$(K,E,N)=(\mathsf{T},\mathsf{U},\mathsf{F})$的合取为假，记录应为未触发、动作为空。只有没有假值时，未知项才使最终结果无法确定。

这条计算规则也适用于前提检查。适用前提未知但触发条件已为假时，$D_v(x)$仍为假，可以确认未触发；适用前提已为假时，则优先按范围外处理。记录要区分已经求值、仍未知和被短路跳过的条件。若否定例外已确认成立，正向规则中的$N$为假，即使一般词语模式命中也应记未触发，并同时保留命中位置和否定依据。

得到候选组合后，需要检查每项条件究竟起了什么作用，而不只是看开发样本能否全部通过。逐项删除条件，可以观察哪些新样本进入触发范围；检查两个条件是否总是同时出现，则有助于找出可能的冗余。对合取规则增加条件通常会缩小覆盖，但新增条件本身也可能错误，不能由覆盖变窄推出可靠性提高。从数据归纳或由模型生成的组合同样需要这些检查；组合越深，需要调试的路径和前处理依赖也越多。

列表与树的控制策略规定多条规则项或条件分支怎样选择输出，合取短路则规定同一条件组合内部哪些计算可以跳过。它们与分别保留输出的独立规则来源不同：多条独立规则给出不同标签时，若只按执行顺序覆盖或丢弃输出，就会隐藏冲突，不能说明哪个标签有依据。应先保留各条规则的版本、证据和原始输出，再交由第~\ref{chap:knowledge-fusion}章的多来源融合过程处理。

\section{基于知识关联的规则利用}
\label{sec:rule-based-knowledge-acquisition-association}

在任命新闻中，即使已经找到“林舟”，这个名字也不能说明他是哪一位同名作者；即使确认了“首席架构师”，还要知道任务采用的岗位体系怎样分类。这里缺少的不是另一个字面模式，而是样本与已有知识之间的对应。\term{知识关联}（knowledge association）先确定样本指向哪个已有实体、概念或事实，再用这一对应产生候选知识。

查到一项知识以后，仍有两个问题要分别回答：这项知识本身是否成立，它能否用于当前样本。实体属性可能已经过时，岗位体系可能服务于另一种分类目的，真实的知识库关系也未必在当前句中表达。明确记录概念及其关系的体系称为\term{本体}（ontology），查询它时需要核对采用的分类口径。表~\ref{tab:rule-based-knowledge-acquisition-association}按实体、概念和关系三种查询对象，比较需要建立的对应及错误位置。同一任务可以使用多种关联，但不必固定按三者依次执行。

\begin{table*}[htbp]
  \centering
  \begingroup
  \setlength{\tabcolsep}{2.5pt}
  \setlength{\parindent}{0pt}
  \setlength{\RaggedRightParindent}{0pt}
  \renewcommand{\arraystretch}{1.12}
  \begin{tabularx}{\linewidth}{>{\raggedright\arraybackslash}p{17mm}>{\raggedright\arraybackslash}p{25mm}>{\raggedright\arraybackslash}p{26mm}>{\raggedright\arraybackslash}p{30mm}>{\raggedright\arraybackslash}p{29mm}>{\raggedright\arraybackslash}X}
    \toprule
    关联类型 & 样本线索 & 查询对象 & 对应条件 & 候选知识 & 主要歧义 \\
    \midrule
    实体匹配 & 名称、别名、标识及上下文 & 实体标识及属性 & 提及指向该实体；属性适合当前时间与任务 & 类别、属性或查询入口 & 同名未必同实体；无候选与多候选有别 \\
    概念关联 & 词语、短语及局部用法 & 词典、本体或分类概念 & 表达指向该概念；目标属性允许沿层级继承 & 概念标签、上位类或候选范围 & 相似未必同义；上位类不能确定唯一细类 \\
    关系匹配 & 已消歧实体对及证据范围 & 有方向和时间的关系事实 & 实体与关系定义对齐；句级需表达证据，组级需另立假设 & 句级关系候选或提及组监督 & 共现未必表达；缺失未必为假 \\
    \bottomrule
  \end{tabularx}
  \endgroup
  \caption{三类知识关联的输入、依据与错误位置。表中每一行都要求同时保存样本侧证据和知识侧版本，避免只留下最终标签而无法判断错误发生在哪一步。}
  \label{tab:rule-based-knowledge-acquisition-association}
\end{table*}

\subsection{基于实体匹配的规则利用}
\label{sec:rule-based-knowledge-acquisition-entity-matching}

要查询新闻中“林舟”的属性，我们先要确定这个名字指的是谁。本章用\term{实体匹配}（entity matching）表示样本对象与已有实体之间的对应；文本中的具体操作称为\term{实体链接}（entity linking），即把一个提及对应到知识库中有唯一标识的对象。规则由此把检查对象从文字跨度转为可查询的实体标识，而不是只比较字符串。输入包括提及跨度、名称、别名、标识符和上下文；输出可以是唯一实体、多个候选、没有候选，或在证据不足时弃权。

设教学实体库有两位名为“林舟”的作者：\texttt{E-017}的记录涉及海岬科技，\texttt{E-204}的记录涉及东岭大学。对“林舟介绍了海岬科技的新系统”，规则先按名字生成两个候选，再结合句中组织、作者主页标识或文档来源消歧，也就是判断提及应当对应哪位作者。句中出现海岬科技是一条待检查的线索，名字相同本身并没有确定身份。如果只有“林舟发表了报告”，应保留$\{\texttt{E-017},\texttt{E-204}\}$而不是任取一位。库中没有候选时，需要补充实体或检索线索；有多个候选无法区分时，需要补充区分身份的证据，两种情况应分别保存。

若提及带有确属该对象的唯一标识，按标识查询可直接确定实体；别名表则把简称和旧名作为查询入口，需要记录名称有效的时间。模糊字符串匹配放宽了拼写要求，既能找到更多变体，也可能把不同名字混在一起。模型可以利用更广的上下文辅助消歧，承担的仍是关联对象的选择。实体链接研究把候选生成与基于上下文的候选选择分开，有助于区分“没有检索到实体”和“在候选中选错实体”\scite{knowledge-acquisition-milne2008}

确定实体以后，才能查询属性，但查到的属性仍要与标注任务对照。例如，\texttt{E-017}的职业记录为“软件工程师”，任务要求的却是文档所述时间的职务，静态职业不能直接替代当时的任职信息。因此要区分事实本身有误、提及链接错误和属性不适用于当前对象三类问题，并保存实体库版本、属性生效时间、候选列表、消歧证据与最终选择。

维护实体词表时，专家可以补充别名，也可以从历史匹配数据归纳或采纳模型建议。新增别名仍需检查是否混入不同实体的名字、是否只在特定时间有效，并通过旧记录确认改动影响哪些结果。若实体间可用别名和上下文记录的数量或质量不同，需要按目标实体组分别检查候选召回与选择错误：前者检查正确实体是否进入候选，后者检查候选中是否选错身份。材料多不保证召回更高，也不保证消歧正确；资料充分的实体上的结果不能代表资料不足的对象。

\subsection{基于概念关联的规则利用}
\label{sec:rule-based-knowledge-acquisition-concept-association}

新闻中的“首席架构师”可能需要转换为任务要求的岗位类别，这与确定“林舟”是哪位作者是不同问题。这里查询的是岗位含义，不以先确定人员身份为必要条件。\term{概念关联}（concept association）把词语或短语对应到词典、本体或分类体系中的概念标识，再按体系明确的上下位关系产生类别或候选范围。文字用于找到概念，推导则依据概念标识与体系中的关系。

假设教学岗位体系把“首席架构师”对应到概念\texttt{role:\allowbreak{}technical-\allowbreak{}lead}，并把该概念归入上位类\texttt{role:\allowbreak{}management-\allowbreak{}position}。文本确认细类后，规则可以沿“属于”关系给出“管理岗位”这一上位类别。只知道管理岗位时却不能反推首席架构师，因为该上位类还包含其他岗位。粗类别任务可以使用上位概念；细类别任务若没有进一步证据，只能保留下位候选集合。这里的归类依据是指定教学体系，不是对所有岗位体系的断言。

建立词语到概念的对应时，还要区分词语之间是什么关系。在当前用法下的同义词可以指向同一概念，缩写需要结合领域与上下文确定展开含义；近义词仅表示含义相近，不足以认定为同一概念，上下位词则不能合并掉层级方向。WordNet把同义词组织为同义词集，再用语义关系连接概念，提供了分开记录词面集合与概念关系的实例\scite{knowledge-acquisition-miller1995}

为了追查归类依据，概念映射表、本体关系和推导规则需要分别维护。映射表确定文本指向哪个概念，本体给出概念间的关系，推导规则规定哪些关系可以转成任务标签。映射候选可由人工整理、词语使用模式归纳或模型建议得到，但都要分别检查表达含义与体系关系。属性的继承也需核对语义；细类的成员属于管理岗位，不意味着某个具体岗位的薪酬、任期或资格要求适用于全部父子概念。

同一个岗位词在按职责、组织级别或行业标准建立的体系中，可能对应不同分类位置。规则因此要保存采用的体系与版本，概念合并、拆分或层级调整时，沿映射和推导记录找到需要重算的旧结果。当前体系缺少映射只说明规则没有找到关联，不能证明样本不属于目标类别。跨来源概念怎样对齐，由第~\ref{chap:knowledge-fusion}章继续讨论。

\subsection{基于关系匹配的规则利用}
\label{sec:rule-based-knowledge-acquisition-relation-matching}

若任务要标注“林舟任职于海岬科技”，查询单个实体属性或岗位类别还不够，需要检查人员与组织之间的关系。\term{关系匹配}（relation matching）把样本中的实体对与知识库关系对应起来，用关联事实产生关系标签或成组监督。检查对象由单个实体转为有次序的实体对，关系定义及有效时间也要随之确定。例如，$(\text{林舟},\text{任职于},\text{海岬科技})$不能换成反向三元组，也不能用于事实生效前后的所有文本。

没有逐句人工关系标签时，可以先利用知识库事实寻找共同出现相关实体的文本，构造训练候选。这种借助已有事实提供标签的办法称为\term{远程监督}（distant supervision）。Mintz等人的设定将知识库实体关系与文本中共同出现的实体对对齐，以此构造训练实例\scite{knowledge-acquisition-mintz2009} 对齐提供了启发式监督，但没有逐句证明文本表达了那项事实。

我们继续使用构造事实“林舟任职于海岬科技”，比较两句话。本例把目标关系限定为文本所述时间已经生效的任职状态，并假设句甲中的任命已经生效、该时点仍在知识库事实的有效期内。在这一教学前提下，句甲“海岬科技任命林舟为首席架构师”支持目标任职关系，句乙“林舟在论坛上批评了海岬科技的发布流程”只同时提及两者。把事实复制给所有共现句子，就会给句乙添加没有文本依据的正标签。这是教学中已知的语义对照，并不是机器仅凭共现得到的判断。图~\ref{fig:rule-based-knowledge-acquisition-relation-path}中，两句都经过实体消歧和关系对齐后归入提及组；表示语义对照的点线不属于自动判定步骤，构造组也不要求预先知道哪句表达了关系。

\begin{figure*}[htbp]
  \centering
  \begin{tikzpicture}[
  x=1cm,
  y=1cm,
  font=\figurefont\footnotesize,
  source/.style={draw=FigureInput,line width=0.9pt,fill=FigureInput!9,minimum height=10mm,text width=32mm,align=left,inner sep=1.5mm},
  fact/.style={draw=FigureModel,line width=1pt,fill=FigureModel!10,minimum height=9mm,text width=53mm,align=center,inner sep=1.5mm},
  process/.style={draw=BookMuted,line width=0.8pt,fill=BookPaper,minimum height=10mm,text width=13mm,align=center,inner sep=1.5mm},
  comparison/.style={draw=BookMuted,line width=0.8pt,dotted,fill=white,minimum height=15mm,text width=36mm,align=center,inner sep=1.5mm},
  group/.style={draw=FigureOutput,line width=1.1pt,fill=FigureOutput!7,minimum height=37mm,text width=36mm,align=center,inner sep=1.5mm},
  flow/.style={knowledge arrow,draw=BookInk,line width=1.1pt,rounded corners=1mm},
  evidence/.style={knowledge arrow,draw=FigureModel,line width=0.9pt,dashed,rounded corners=1mm},
  contrast/.style={draw=BookMuted,line width=0.8pt,dotted},
  edge label/.style={text=BookMuted,inner sep=1pt,font=\figurefont\scriptsize}
]
  \fill[white] (0,0) rectangle (16.7,7);

  \node[source] (sent-a) at (1.8,5.35) {句甲：海岬科技任命林舟为首席架构师。};
  \node[source] (sent-b) at (1.8,2.25) {句乙：林舟在论坛上批评了海岬科技的发布流程。};

  \node[process] (entity-a) at (4.85,5.35) {实体消歧};
  \node[process] (entity-b) at (4.85,2.25) {实体消歧};
  \node[process] (mapping-a) at (7.1,5.35) {关系对齐};
  \node[process] (mapping-b) at (7.1,2.25) {关系对齐};

  \node[fact] (kb) at (7.35,3.8) {已有构造事实\\$(\text{林舟/E-017},\text{任职于},\text{海岬科技/E-003})$};
  \node[comparison] (positive) at (10.3,5.35) {教学语义对照\\句甲表达关系\\可作句子级候选};
  \node[comparison] (cooccur) at (10.3,2.25) {教学语义对照\\句乙仅共同提及\\不应给句子正标签};
  \node[group] (bag) at (14.55,3.8) {提及组级候选\\句甲、句乙组成一组\\由实体关系库事实关联关系\\假设至少一条表达关系};

  \draw[flow] (sent-a) -- (entity-a);
  \draw[flow] (entity-a) -- (mapping-a);
  \draw[contrast] (mapping-a.east) -- (positive.west);
  \draw[flow] (mapping-a.north) -- (7.1,6.55)
    -- node[edge label,above] {按实体对归组} (14.55,6.55) -- (bag.north);

  \draw[flow] (sent-b) -- (entity-b);
  \draw[flow] (entity-b) -- (mapping-b);
  \draw[contrast] (mapping-b.east) -- (cooccur.west);
  \draw[flow] (mapping-b.south) -- (7.1,1.05)
    -- (14.55,1.05) -- (bag.south);

  \draw[evidence] (kb.north -| mapping-a.south) -- (mapping-a.south);
  \draw[evidence] (kb.south -| mapping-b.north) -- (mapping-b.north);
  \draw[evidence] (kb.east) -- (bag.west);

  \draw[flow] (0.35,0.35) -- (1.25,0.35);
  \node[anchor=west,text=BookMuted,font=\figurefont\scriptsize] at (1.35,0.35) {处理步骤};
  \draw[evidence] (4.1,0.35) -- (5.0,0.35);
  \node[anchor=west,text=BookMuted,font=\figurefont\scriptsize] at (5.1,0.35) {事实关联};
  \draw[contrast] (7.85,0.35) -- (8.75,0.35);
  \node[anchor=west,text=BookMuted,font=\figurefont\scriptsize] at (8.85,0.35) {已知语义对照，非自动判定};
\end{tikzpicture}

  \caption{从已有事实到文本候选关系的关联路径。实体和关系均为构造例子，并假设句甲的任命已生效、所述时点在事实有效期内；实线箭头表示处理步骤，虚线箭头表示知识库事实的证据关联，点线引出教学中已知的语义对照，不表示自动判定。实体消歧确认两句中的提及指向同一实体对；关系对齐检查定义、方向与时间。在上述前提下，句甲支持任职关系，句乙仅共同提及，但两句都可按实体对进入同一提及组。组级候选依赖至少一次表达的假设，不要求预先知道哪句支持关系，也不能据此给每句确定的正标签。}
  \label{fig:rule-based-knowledge-acquisition-relation-path}
\end{figure*}

这一路径要先确认实体及关系方向，再核对知识库关系与任务关系的定义，最后规定证据应在当前句、整段还是同一实体对的多条提及中出现。证据范围改变，监督对象也随之改变。Riedel等人分别建模“实体对存在关系”和“某个提及表达该关系”，采用至少一次表达的假设，放松了每个共同提及都必须表达关系的要求\scite{knowledge-acquisition-riedel2010}

采用这一假设时，我们把同一实体对的多条提及作为一个组，用组级候选表示“组内至少一条可能支持关系”，而不把每句都标为正例。训练过程可以先为各句是否表达关系评分，再聚合为组级预测，只与组级标签比较。第~\ref{chap:knowledge-acquisition-model-based}章的式~\eqref{eq:knowledge-acquisition-model-bag-prediction}给出计算与假设。本节规则提供的是提及组、实体标识和关联事实，具体哪句构成证据仍需提及级推断，不能拿教学中已知的对照代替机器输出。

保存这条路径，也能区分不同错误：消歧错误关联了无关实体，过时事实把关系用于错误时间，文本未表达关系则未满足指定证据范围。另一种情况是知识库根本没有记录该关系。这里采用\term{开放世界假设}（open-world assumption），即知识库可能不完整，“未查到”不等于“已证明不存在”。规则应保留未知或候选关系，只有明确的排除证据符合负类定义时才输出负标签。

即使上述步骤都能执行，仍需分别检查各关系类别的候选总量与每个实体对的提及组大小。类别总量不同，可能来自每类包含的实体对数不同，不能由此推出单个实体对的提及更多。若同一实体对反复共现，或其文本被重复转载，在未去重时便会增加该对的提及数。仅看执行是否成功无法发现这些分布差异，因而还要保存知识库版本、事实来源、实体对、提及列表和去重标识，供后续融合与可靠性检查追查。

\section{基于约束推理的规则利用}
\label{sec:rule-based-knowledge-acquisition-constraint}

对文本中的同一岗位分别作出细类和上位类判断后，仍可能出现互相矛盾的组合。例如，细类被判为真，其必属的上位类却被判为假。这里可用的依据不是再分别指定答案，而是规定这些答案必须满足什么关系。\term{约束推理}（constraint-based inference）把未知标签的可能取值列为候选，再用逻辑、结构或领域条件排除不允许的组合；仍有多个候选时，可以保留范围，或按允许违反但需付出代价的软偏好选择。

要判断约束究竟排除了什么，我们需要先确定任务允许哪些完整答案，再计算其中哪些满足全部必须遵守的条件。设全部候选答案为$\bm{Y}$，其中每个$\bm{y}$是一组未知标签或数值的完整赋值；任务规定的取值范围确定了$\bm{Y}$。硬约束$c_1,\ldots,c_m$表示必须满足的条件，可行集合为
\begin{equation}
  \bm{F}
  =
  \{\bm{y}\in\bm{Y}:c_k(\bm{y})=\mathsf{T},\ k=1,\ldots,m\}.
  \label{eq:rule-based-knowledge-acquisition-feasible-set}
\end{equation}
式中要求每一条硬约束都为真，因此$\bm{F}$回答的是哪些完整答案仍被允许，而不是哪个答案更受偏好。在约束能够对候选求值的前提下，它可能只有一个元素，也可能有多个元素或为空。唯一解表示当前前提恰好确定候选，多解表示依据只缩小了范围，空集则提示已知事实、约束或输入之间存在冲突。缺少前提而无法判断约束时，应保留无法确定，不能直接宣告无解。求解应分别报告这些状态，不能在多解时任取一项当作已确定知识，也不能为得到输出而悄悄放松约束。条件检查触发了求解，不等于求解已经给出任务内容。任务允许返回范围时，非空可行集可以作为内容提交，来源动作是作答、语义为已给出，但单值仍可能未知；无法判断或已证明无解时，本节都规定不提交任务赋值，动作是弃权、内容为空、语义为未知。前者保存缺少的前提，后者另存无解诊断和冲突依据，不能把二者混为一类。上述执行状态均以相应检查或求解已完成为成功条件；程序中断则记失败，验证状态也不因求解完成而自动变为已核验。

\subsection{基于逻辑约束的规则利用}
\label{sec:rule-based-knowledge-acquisition-logical-constraint}

要排除“细类成立、上位类不成立”的组合，可以把两项类别判断写成命题，再规定它们之间的蕴含。\term{逻辑约束}（logical constraint）用蕴含、互斥或等价等关系，限制命题可以怎样共同取值。作为硬约束时，违反关系的候选被排除；作为\term{软约束}（soft constraint）时，违反仍被允许，但要付出指定代价。后者表达的是偏好，需要结合权重或其他比较依据来选择答案。

设二值命题$p$表示“属于技术管理岗位”，$q$表示“属于管理岗位”，教学体系要求$p\Rightarrow q$。两个命题共有四种赋值，其中$(p,q)=(1,0)$违反蕴含，被硬约束排除，另外三种仍可行。若再人为设置软约束代价$L(\bm{y})$，就可以按代价从小到大比较可行候选。图~\ref{fig:rule-based-knowledge-acquisition-feasible-set}给出教学代价：三个可行项中最小者唯一，但唯一最小代价只说明它符合当前偏好，没有证明它是真实答案。下方面板则显示，加入相互矛盾的已知条件会使可行集变空。

\begin{figure*}[htbp]
  \centering
  \begin{tikzpicture}[
  x=1cm,
  y=1cm,
  font=\figurefont\footnotesize,
  candidate/.style={draw=BookMuted,line width=0.8pt,fill=BookPaper,minimum width=30mm,minimum height=13mm,align=center,inner sep=1.5mm},
  selected/.style={candidate,draw=FigureModel,line width=1.2pt,fill=FigureModel!10},
  premise/.style={draw=FigureInput,line width=0.9pt,fill=FigureInput!9,minimum width=28mm,minimum height=10mm,align=center,inner sep=1.5mm},
  constraint/.style={draw=BookAmber,line width=1pt,fill=BookAmber!9,minimum width=31mm,minimum height=10mm,align=center,inner sep=1.5mm},
  combine/.style={circle,draw=BookInk,line width=0.9pt,fill=white,minimum size=7mm,inner sep=0pt},
  empty/.style={draw=BookAmber,line width=1.2pt,fill=BookAmber!9,minimum width=32mm,minimum height=11mm,align=center},
  flow/.style={knowledge arrow,draw=BookInk,line width=1.1pt,rounded corners=1mm},
  soft/.style={text=BookMuted,font=\figurefont\scriptsize}
]
  \fill[white] (0,-1.15) rectangle (16.6,6.5);
  \node[anchor=north west,font=\figurefont\bfseries] at (0,6.25) {(a) 硬约束筛选与软偏好};
  \node[constraint,minimum width=44mm] (hard) at (8.3,5.75) {硬约束：$p\Rightarrow q$};

  \node[candidate] (c00) at (1.8,4.25) {$(p,q)=(0,0)$\\{\scriptsize 可行；$L=1.0$}};
  \node[candidate] (c01) at (5.65,4.25) {$(p,q)=(0,1)$\\{\scriptsize 可行；$L=0.6$}};
  \node[candidate,draw=BookAmber,fill=BookAmber!7] (c10) at (9.5,4.25) {$(p,q)=(1,0)$\\{\scriptsize 违反硬约束}};
  \node[selected] (c11) at (13.35,4.25) {$(p,q)=(1,1)$\\{\scriptsize 可行；$L=0.2$（最小）}};
  \draw[BookAmber,line width=1.2pt] ([xshift=-10mm,yshift=-4mm]c10.center) -- ([xshift=10mm,yshift=4mm]c10.center);
  \draw[BookAmber,line width=1.2pt] ([xshift=-10mm,yshift=4mm]c10.center) -- ([xshift=10mm,yshift=-4mm]c10.center);
  \node[soft,anchor=north] at (8.3,3.42) {硬约束留下三个答案；软代价只负责排序};

  \draw[BookRule,line width=0.8pt] (0,2.95) -- (16.6,2.95);
  \node[anchor=north west,font=\figurefont\bfseries] at (0,2.75) {(b) 新增互相矛盾的已知条件};
  \node[premise] (p-one) at (2.2,1.65) {观测：$p=1$};
  \node[premise] (q-zero) at (7.2,1.65) {观测：$q=0$};
  \node[constraint] (imply) at (12.2,1.65) {硬约束：$p\Rightarrow q$};
  \node[combine] (and) at (7.2,-0.35) {$\land$};
  \node[empty] (none) at (12.4,-0.35) {无可行解\\报告冲突};
  \draw[flow] (p-one.south) |- (and.west);
  \draw[flow] (q-zero.south) -- (and.north);
  \draw[flow] (imply.south) -- ++(0,-0.5) -| (and.45);
  \draw[flow] (and.east) -- (none.west);
\end{tikzpicture}

  \caption{约束缩小可行答案并产生冲突的教学示意。上面板从两个二值命题的四种赋值出发，硬约束$p\Rightarrow q$排除$(1,0)$，软约束代价再在三个可行候选中给出唯一最小值；这些代价是人为设置，不是统计估计。下面板同时加入观测$p=1$、观测$q=0$和同一蕴含后，可行集合变为空集。逻辑可行、偏好最小和真实正确是三种不同判断。}
  \label{fig:rule-based-knowledge-acquisition-feasible-set}
\end{figure*}

当多条软约束分别表达偏好时，我们需要一个可计算的量比较它们共同作用下的候选。令第$j$条软约束的违反代价为$\ell_j(\bm{y})\geq0$，零表示不违反，正值表示违反程度；对应权重为$w_j\geq0$，用于规定该项代价在比较中的分量。候选仍限于已经满足全部硬约束的$\bm{F}$，对选定的有限条软约束将加权代价相加，数值越小表示越受偏好。可行集非空且最小值能够取到时，按这个总量选择候选，得到
\begin{equation}
  \widehat{\bm{y}}
  \in
  \argmin_{\bm{y}\in\bm{F}}
  \sum_j w_j\ell_j(\bm{y}).
  \label{eq:rule-based-knowledge-acquisition-soft-constraint}
\end{equation}
这里$\bm{y}\in\bm{F}$保证候选没有违反硬约束，目标值较小才表示按当前权重更受偏好。$\argmin$保留所有达到最小值的候选；若最小值并列，就仍未唯一确定答案。有限非空候选集上的有限代价会有最小值，连续取值问题则需另外确认前述存在条件。约束推断可以通过整数或线性规划整合局部预测与任务条件，Roth和Yih给出了自然语言任务中的代表性形式\scite{knowledge-acquisition-roth2004} 若要把加权逻辑进一步解释为概率模型，还需确定权重的统计含义，以及公式替换为具体对象后的实例数量；Markov逻辑网络提供了这种概率化实现\scite{pgm-energy-models-richardson2006}

写出逻辑约束时，要保留原关系的方向与适用对象。由“首席架构师属于技术负责人”只能得到$\text{首席架构师}\Rightarrow\text{技术负责人}$，不能反推所有技术负责人都是首席架构师；针对正式岗位的互斥条件也未必适用于临时角色。专家给出的约束、从数据归纳的约束和模型提出的候选，都需要核对方向、量词、对象与例外。

为了在约束更新后重新检查结果，推理记录要保存使用的前提、被排除的候选，以及各条软约束贡献的代价。求解所需计算、软权重设定和冲突解释也都计入成本。这里讨论的是从前提产生或限制新候选；约束若主要用于综合多个现成来源，则由第~\ref{chap:knowledge-fusion}章的融合方法处理。

\subsection{基于结构约束的规则利用}
\label{sec:rule-based-knowledge-acquisition-structural-constraint}

文本中的实体片段可能跨越多个词元，逐个位置选择标签后，还要检查它们能否拼成合法片段。\term{结构约束}（structural constraint）利用序列、树、图或成组记录中的组织关系，限制标签组合。与一般命题关系相比，这里还要明确哪些位置或对象相邻、包含或属于同一组；这些结构可由任务格式、标注规范或输入给出，也可由实体匹配等步骤构造。

为了让逐词标签表达片段边界，本节采用BIO约定，其中B、I、O分别来自beginning、inside和outside：\texttt{B-X}开始类型X的片段，\texttt{I-X}延续同类型片段，\texttt{O}位于片段外。因此，\texttt{I-X}只能跟在\texttt{B-X}或\texttt{I-X}之后，不能凭空开始一个片段。早期文本分块工作把片段结构编码成逐词标签，提供了按边界表示约定检查局部标签的实例\scite{knowledge-acquisition-ramshaw1995}

对教学序列“林舟 加入 海岬 科技”，先由局部规则为每个位置给出候选，标签中的\texttt{PER}与\texttt{ORG}分别表示人员和组织。若逐点选择得到\texttt{B-PER, I-ORG, B-ORG, I-ORG}，第二个位置想延续组织片段，前一个位置却开始了人员片段，因此组合非法。图~\ref{fig:rule-based-knowledge-acquisition-sequence-constraint}中的\texttt{B-PER, O, B-ORG, I-ORG}满足转移条件。然而，把“海岬科技”全部标为\texttt{O}同样合法，却漏掉了组织提及。结构检查排除非法衔接后，仍需词语、上下文或其他来源的语义评分，才能选择任务答案。

\begin{figure*}[htbp]
  \centering
  \begin{tikzpicture}[
  x=1cm,
  y=1cm,
  font=\figurefont\footnotesize,
  token/.style={draw=FigureInput,line width=0.9pt,fill=FigureInput!9,minimum width=30mm,minimum height=8mm,align=center},
  tag/.style={draw=FigureModel,line width=0.9pt,fill=FigureModel!9,minimum width=30mm,minimum height=8mm,align=center},
  bad tag/.style={tag,draw=BookAmber,fill=BookAmber!9},
  choices/.style={draw=BookMuted,line width=0.7pt,dashed,fill=white,minimum width=34mm,minimum height=9mm,align=center,font=\figurefont\scriptsize},
  allowed/.style={knowledge arrow,draw=BookInk,line width=1pt},
  invalid/.style={knowledge arrow,draw=BookAmber,line width=1pt,dashed},
  edge label/.style={fill=white,text=BookMuted,inner sep=1pt,font=\figurefont\scriptsize}
]
  \fill[white] (0,0) rectangle (16.85,6.3);
  \node[anchor=east,font=\figurefont\bfseries] at (1.25,5.85) {词元};
  \foreach \x/\word in {3.1/林舟,7.3/加入,11.5/海岬,15.1/科技} {
    \node[token] at (\x,5.85) {\word};
  }

  \node[anchor=east,font=\figurefont\bfseries,text=BookAmber] at (1.25,4.45) {局部};
  \node[tag] (bad1) at (3.1,4.45) {\texttt{B-PER}};
  \node[bad tag] (bad2) at (7.3,4.45) {\texttt{I-ORG}};
  \node[tag] (bad3) at (11.5,4.45) {\texttt{B-ORG}};
  \node[tag] (bad4) at (15.1,4.45) {\texttt{I-ORG}};
  \draw[invalid] (bad1) -- (bad2);
  \draw[allowed] (bad2) -- (bad3);
  \draw[allowed] (bad3) -- (bad4);
  \node[circle,draw=BookAmber,fill=white,text=BookAmber,line width=0.8pt,minimum size=4mm,inner sep=0pt,font=\figurefont\bfseries\scriptsize] at (5.2,4.45) {$\times$};

  \node[anchor=east,font=\figurefont\bfseries,text=FigureModel] at (1.25,3.0) {合法};
  \node[tag] (good1) at (3.1,3.0) {\texttt{B-PER}};
  \node[tag] (good2) at (7.3,3.0) {\texttt{O}};
  \node[tag] (good3) at (11.5,3.0) {\texttt{B-ORG}};
  \node[tag] (good4) at (15.1,3.0) {\texttt{I-ORG}};
  \draw[allowed] (good1) -- node[edge label,above] {允许} (good2);
  \draw[allowed] (good2) -- node[edge label,above] {允许} (good3);
  \draw[allowed] (good3) -- node[edge label,above] {允许} (good4);

  \node[anchor=east,font=\figurefont\bfseries,text=BookMuted] at (1.25,1.45) {候选};
  \node[choices] at (3.1,1.45) {\{\texttt{B-PER}, \texttt{O}\}};
  \node[choices] at (7.3,1.45) {\{\texttt{O}, \texttt{I-ORG}\}};
  \node[choices] at (11.5,1.45) {\{\texttt{B-ORG}, \texttt{O}\}};
  \node[choices] at (15.1,1.45) {\{\texttt{I-ORG}, \texttt{O}\}};
  \node[anchor=south,text=BookMuted,font=\figurefont\scriptsize,align=center] at (8.4,0.2) {结构约束排除非法组合，但合法序列仍可能不止一个};
\end{tikzpicture}

  \caption{局部标签与整体结构的约束关系。采用BIO边界约定：\texttt{B-X}开始类型X片段，\texttt{I-X}只能延续同类型的\texttt{B-X}或\texttt{I-X}，\texttt{O}位于片段外。上行的\texttt{B-PER}\,$\to$\,\texttt{I-ORG}违反转移条件；下行是一种合法候选。每个位置下方保留局部候选集合，联合约束筛选后仍可能留下多个合法序列。}
  \label{fig:rule-based-knowledge-acquisition-sequence-constraint}
\end{figure*}

对这个序列，可以在相邻位置间反复删除找不到合法相邻标签的候选，这就是局部传播。图中第二位置的\texttt{I-ORG}在第一位置的$\{\texttt{B-PER},\texttt{O}\}$中找不到同类型前驱，因此先删除它；反复检查到不再删除候选时停止。第四位置的\texttt{I-ORG}仍有第三位置的\texttt{B-ORG}支持，不能删除，但第三位置选\texttt{O}时，第四位置就不能选\texttt{I-ORG}。所以，传播后各位置仍有候选，不意味着任意组合都是合法序列。若要给出完整赋值，或按语义评分选择最优合法序列，需要继续作联合选择，可按结构使用动态规划、搜索或约束求解；一般图上的局部支持还不能保证全局有解。扩展到边界包含、父子标签或跨记录身份一致性时，要另行说明检查的结构关系。尤其是同一实体跨记录应保持标识一致，并不要求它的岗位、位置或状态一直不变。

这些推理都依赖输入结构有依据。文档层级可能直接来自文件格式，跨记录的同一对象可能由实体匹配确定，图中邻接边可能由距离或关系规则构造。缺边会阻断应有的传播，错边会把无关对象联结起来，错误边界则可能把本来合法的标签判为冲突。因此，结构的来源与版本要随标签证据保存，求解成功不能反过来证明结构正确。

求解成本还取决于结构大小及连接方式。只含局部转移条件的长序列可以利用动态规划，稠密图中的全局条件却可能要求更大的搜索空间。使用近似方法时，记录应说明尚未解决的冲突，以及输出究竟是完整解还是仅满足部分局部条件的结果。

\subsection{基于领域约束的规则利用}
\label{sec:rule-based-knowledge-acquisition-domain-constraint}

盘点物料时，初始量、流入量和结束量可以限制尚未记录的流出量；但这要以没有其他来源或损失为前提。\term{领域约束}（domain constraint）把守恒、数量平衡、范围限制或业务一致性关系写成未知标签或数值必须满足的条件，再用观测缩小候选。这里改变的是约束依据：它来自特定领域，不一定需要换求解器，但必须说明变量含义、适用环境、单位与容差。

为计算流出量，我们构造同一核算时段内的教学观测：初始量$s_0=100\,\mathrm{kg}$，流入量$q_{\mathrm{in}}=30\,\mathrm{kg}$，结束量$s_1=108\,\mathrm{kg}$，未知量为$q_{\mathrm{out}}$。没有其他流入、流出或损失时，结束量应等于初始量加流入量再减流出量，因此使用精确关系
\begin{equation}
  s_1=s_0+q_{\mathrm{in}}-q_{\mathrm{out}},
  \label{eq:rule-based-knowledge-acquisition-balance}
\end{equation}
可求得$q_{\mathrm{out}}=22\,\mathrm{kg}$。等式把三个已知量与一个未知量联系起来，完整观测和精确平衡前提共同确定了这个唯一值。若还有未记录的来源或损失，就不能沿用这条等式作同样推断。

测量有误差时，需要改为判断哪些流出量仍与观测相容。这里用整体平衡残差，也就是观测结束量与平衡关系计算值之差，量化偏离；取绝对值是为了同时限制两个方向的偏离。假设各观测误差合成后的整体残差绝对值不超过$2\,\mathrm{kg}$，则只要求
\[
  \left|s_1-(s_0+q_{\mathrm{in}}-q_{\mathrm{out}})\right|
  \leq 2\,\mathrm{kg},
\]
代入同一组数值，规则只能推出$q_{\mathrm{out}}\in[20,24]\,\mathrm{kg}$，不能把区间中心当作已经确定的答案。这里的$2\,\mathrm{kg}$约束整体残差，不是三个观测各自的误差界。流入量缺失时，原观测不再足以确定这一范围；若结束量改为$s_1=140\,\mathrm{kg}$且仍要求$q_{\mathrm{out}}\geq0$，在无其他来源的前提下则无解。记录要分别返回唯一值、候选范围或观测冲突，并保存单位、容差与求解精度。

是否可以要求精确平衡，要看这条关系的依据。按任务定义成立的业务恒等式可以作为硬约束；领域规律若依赖封闭、稳态或理想条件，就要把这些条件写入适用前提。从数据观察到的经验关系通常更适合写成软约束，并保留估计误差与更新条件。表~\ref{tab:rule-based-knowledge-acquisition-domain-constraint}固定物料平衡例子，比较三种依据分别允许得到什么结论，不能只因为表达式相似就互换它们。

\begin{table*}[htbp]
  \centering
  \begingroup
  \setlength{\tabcolsep}{2.5pt}
  \setlength{\parindent}{0pt}
  \setlength{\RaggedRightParindent}{0pt}
  \renewcommand{\arraystretch}{1.12}
  \begin{tabularx}{\linewidth}{>{\raggedright\arraybackslash}p{23mm}>{\raggedright\arraybackslash}p{32mm}>{\raggedright\arraybackslash}p{25mm}>{\raggedright\arraybackslash}p{26mm}>{\raggedright\arraybackslash}p{32mm}>{\raggedright\arraybackslash}X}
    \toprule
    依据类型 & 前提 & 观测 & 误差或容差 & 可得结论 & 失效处理 \\
    \midrule
    按定义成立的关系 & 各量是完整且互斥的记账项 & $s_0=100$、$q_{\mathrm{in}}=30$、$s_1=108$，单位为$\mathrm{kg}$ & 无额外容差 & 精确值$q_{\mathrm{out}}=22\,\mathrm{kg}$ & 报告定义冲突或无解 \\
    理想条件下的领域规律 & 除已列流入、流出外，无其他物质通量或损失；观测误差已合成 & 同上；另核验核算边界及传感器版本 & 整体平衡残差绝对值不超过$2\,\mathrm{kg}$ & 可行范围$q_{\mathrm{out}}\in[20,24]\,\mathrm{kg}$ & 范围外停用；无法确定时弃权 \\
    数据归纳的经验关系 & 关系仅在指定设备、时间与样本范围内观察到 & 同类测量及规则数据版本 & 软代价与经核验的残差范围 & 偏好排序或异常候选 & 条件变化后重新估计 \\
    \bottomrule
  \end{tabularx}
  \endgroup
  \caption{领域约束的前提决定结论强度。表中数值均为单位一致的假设教学数据；精确等式、带容差范围和经验软关系分别支持唯一值、可行区间和偏好候选，不能相互替代。}
  \label{tab:rule-based-knowledge-acquisition-domain-constraint}
\end{table*}

取得候选领域约束后，还要核对变量与数据字段是否表示同一量；模型建议的关系尤其需要检查是否把相关性写成了恒等式。环境改变时，原来没有其他物质通量或损失的前提、误差范围或业务流程可能不再成立，需要重新检查前提。开发和维护因此包括观测对齐、单位换算、误差建模、无解解释与版本管理，而不只是写出公式。

即使平衡求解完成，它也没有独立验证输入测量可靠。按前述状态约定填写本例记录时，假设任务允许返回可行范围，精确例提交$22\,\mathrm{kg}$，带容差例提交$[20,24]\,\mathrm{kg}$；两者均作答、执行成功、语义为已给出、验证为候选，后者没有确定单值。若无法确认是否存在其他物质通量，检查完成后弃权，内容为空、语义为未知，并记录缺少核算边界证据；若确认存在未计入的通量，则该规则不适用，内容与动作为空。若完整观测为$s_1=140\,\mathrm{kg}$且非负流出要求不变，求解成功证明无解，仍不提交流出量，动作是弃权、内容为空、语义为未知；无解诊断及冲突依据另存，不作为流出量标签。无法判断与已证明冲突虽然都没有任务赋值，留下的原因和求解结果却不同；任务允许的多解范围则是可以提交的内容。第~\ref{chap:knowledge-quality-assessment-correction}章将继续检查所得知识的质量并处理错误。

\section{本章小结}
\label{sec:rule-based-knowledge-acquisition-summary}

从新闻中取得任职候选，不能止于程序找到“任命”。我们要逐步确认样本条件、关联对象和答案约束，并说明各步用了什么依据。专家编写、数据归纳和模型生成决定候选规则怎样得到；模式匹配、知识关联和约束推理决定候选知识凭什么产生。这两个问题可以交叉，一条规则也可以使用多种机制，但每一步仍需分别保存输入、输出、证据与版本。

三类依据留下的边界并不相同。模式匹配可以直接输出标签或关系候选，但要维护符号变化、数值基准和组合分支；知识关联可以复用实体、概念与事实，但要检查消歧、版本、时间及证据范围。约束推理排除不合法组合、缩小候选或发现冲突，所得结果仍依赖前提与结构正确。按软约束选出最小代价，只确定了当前偏好下的选择，没有把它变成真实答案。

把任命例子中的步骤连起来：符号模式找到谓词及参与者，实体链接确定“林舟”和“海岬科技”的标识，再按关系定义、方向、时间与证据范围检查任职候选。需要联合决定标签时，还可用已说明的逻辑或结构约束筛选组合。适用前提与触发条件的检查遵循同一三值约定：范围外记不适用；其余情况合取为真才触发，为假则未触发且动作为空，无法确定时弃权。未知项不能遮住已经确认的假值，但列表和树怎样处理未知仍需另定控制策略，后续求解也要单独记录结果。任务允许返回范围时，多个可行解可作为候选内容作答；无解则报告冲突，而不伪造任务标签。没有输出不等于负例，负标签仍需负类定义的明确检查。

规则的明确形式让我们能够沿输入、前提、匹配、关联和求解步骤追查错误，但不会保证所有样本都得到正确答案。维护规则时，要结合可核验依据、适用范围和前处理及更新成本决定是否继续使用。缺少可用条件或知识时补充覆盖，多个来源冲突时进入第~\ref{chap:knowledge-fusion}章，结论是否可靠及怎样修正则由第~\ref{chap:knowledge-quality-assessment-correction}章统一处理。
